\section{Introduction}

The BENEFITS framework \cite{katrakazas2026} proposed a specific answer to a governance problem that had resisted procedural fixes: for migrants, refugees, and low-income participants in social service programs, sovereignty over personal data cannot be established through privacy policies, consent forms, or ethics board approval alone, because these instruments are legible to institutions and opaque to the people they are meant to protect. Its answer was architectural. A community data steward model would place trusted local representatives as mandatory, cryptographically enforced checkpoints between participant data and external release. A consent-filtered synchronization layer would enforce usage control at the device before any record left participant hands. Immutable, participant-readable audit logs would turn an abstract sovereignty claim into something a person without technical training could inspect and understand. Co-design work with 47 migrants and refugees in Greece, and a preliminary evaluation with a subset of 23, gave early support to the underlying hypothesis: participants who could see and verify governance controls trusted the system more than participants working from policy assurances alone.

That evidence, however, was evidence about a design, not about a system. SIDS 2026 is explicit about this in its own account of what had actually been built at the time of writing. Of the three architectural layers it describes, only part was running code: the mobile application, the SQLite and PouchDB local storage with offline queuing, and the CouchDB backend carrying a k-anonymity pipeline. The differential privacy mechanisms, the privacy-budget tracking, and the steward portal itself were, in the paper's own words, still in active development. The steward governance protocol, the mechanism the paper foregrounds as its central governance innovation, quorum-based release authorization, veto rights, hash-chained audit entries, and a CouchDB validation function meant to make consent enforcement unbypassable even by a misconfigured backend, was further back still: described in prose, but by the paper's own account "in the design and specification phase pending pilot deployment." The trust result the paper reports, a rise in self-reported trust score from 2.1 to 3.8 out of 5 across 23 participants, is consequently a measurement of how people responded to a concept and an early interface, not of a mechanism that had been built, adversarially exercised, and shown to hold under load, partition, or attempted tampering. SIDS 2026 does not overstate this. It flags the sample size, the lack of randomization, and the unvalidated single-item trust measure as caveats, and frames its own numbers as suggestive descriptive evidence rather than confirmed causal estimates. The honesty of that framing is precisely what leaves the gap open: a specification, however carefully reasoned, does not verify itself, and a trust score collected against a partially built system cannot speak to whether the finished mechanisms deliver the guarantees claimed for them.

This paper closes that gap. We built the full BENEFITS reference architecture, including the parts SIDS 2026 had left as specification, an EDC-based connector carrying the policy, quorum, privacy, and audit extensions the earlier paper sketched in prose, a working steward-quorum release gate enforcing a signature-count threshold at contract negotiation, a privacy pipeline combining k-anonymity with a Laplace differential-privacy mechanism and an enforcing budget accountant, and a Merkle-epoch audit chain replacing the earlier paper's chained-hash description with a tested, running implementation. Building the steward-quorum gate this way rather than as the cryptographic envelope SIDS 2026 originally described is itself one of the gaps this paper reports rather than glosses over, discussed in full in Section 3. Turning that specification into code did not simply confirm the design. It surfaced questions the specification had no way to anticipate, because they only arise once a mechanism has to actually run. Does a pure-JavaScript Ed25519 verification path inside a CouchDB validation function hold up as batch sizes grow, or does the sovereignty-enforcing S1 filter itself become the bottleneck it was meant to guard against? What epoch length makes a Merkle-batched audit chain both tamper-evident and cheap enough to verify on a participant's device? Does revocation propagate correctly, with no leaked records observed across the simulated partition scenarios tested, when the network partitions, a condition the offline-first design explicitly has to survive rather than assume away? And does the built system enforce the veto guarantee SIDS 2026's own steward-quorum policy specifies, and if not, what would enforcing it cost in release availability? None of these questions could be answered by the design document alone. They required an implementation to interrogate.

The thesis of this paper is that closing the implementation gap and subjecting the resulting mechanisms to direct empirical measurement gives the sovereignty claim a second, independent line of evidence, one that is objective and mechanism-level where SIDS 2026's was subjective and self-reported. We build the full BENEFITS reference architecture and subject each of its sovereignty mechanisms, consent enforcement, steward quorum, audit-chain integrity, and the privacy-utility tradeoff, to direct empirical measurement. What is novel is not any one of these measurement types taken in isolation: quorum-availability tradeoffs and privacy-budget-utility curves, for instance, have real if thin precedent elsewhere in the literature, as Chapter 2 details. What is novel is measuring several of these mechanism classes together, under one built system, and reporting where each actually costs more than the specification implied or degrades under realistic conditions. This is not offered as validation of the earlier trust score. The two are different kinds of evidence answering different questions, and we treat them as complementary legs of a single sovereignty argument rather than as one confirming the other. SIDS 2026 asked whether people would trust an architecture built this way. This paper asks whether the architecture, once built, does what it claims, and at what price. This is also consistent with a broader pattern in systems that formalize privacy guarantees: even at national scale, translating a differential-privacy budget into a usable statistical product carries a real, contested accuracy cost that specification alone cannot resolve \cite{abowd2022}, a tension SIDS 2026 itself anticipated when it warned that its own privacy budget could render aggregate statistics analytically meaningless at the population sizes typical of a single service site. That concern, raised as an open question in the earlier paper, is one of the things this paper is positioned to measure rather than merely state.

The empirical core of the paper is six mechanism-level experiments, each targeting one property the architecture claims and one variable the specification left unresolved. E1 measures the computational cost of Ed25519 signature verification inside the CouchDB consent filter as batch size scales, asking whether the edge-enforcement mechanism itself remains practical at realistic throughput. This measurement runs inside a harness-only design document built for the experiment, not the production S1 filter, which performs no per-record signature verification today. E2 subjects the consent-revocation path to simulated network partition, tracking whether revocation still reaches every replica with no leaked record, under exactly the offline-first conditions the architecture is meant to tolerate. E3 varies the audit chain's epoch size and measures how proof size and inclusion-verification time respond, alongside the chain's ability to detect retroactive tampering. E4 sweeps the differential-privacy budget and measures the resulting accuracy of aggregate statistics, engaging directly with the small-population utility concern SIDS 2026 raised but did not test. G1 compares the built system's veto-blind quorum policy against the veto-enforcing policy SIDS 2026's own text specifies across a swept quorum parameter, measuring the availability cost of enforcing the guarantee the built system currently omits. G2 varies the heterogeneity of participant consent choices and measures its effect on the utility of resulting aggregates, addressing whether the architecture's meet operation over heterogeneous consent produces aggregates too thin to be useful. Each experiment runs on a seeded, reproducible harness with committed output, so that its result can be checked rather than taken on the paper's word.

The remainder of the paper is organized as follows. Section 2 (Related Work) extends SIDS 2026's own positioning one layer further, from governance models that claim particular properties to systems, including this one, whose properties are independently verified. Section 3 (System Architecture and Implementation) describes the built reference architecture in full, including the components that existed only as specification when SIDS 2026 was written. Section 4 (Experimental Methodology) details the design of the six experiments introduced above. Section 5 (Results) reports what each experiment found. Section 6 (Discussion) situates those findings against SIDS 2026's own open governance questions and against the complementary-evidence relationship between the two papers' trust and mechanism-level claims. Section 7 (Conclusion) restates what the six experiments collectively establish about sovereignty-as-architecture. Section 8 (Future Work) sets out what remains before the three-country pilot SIDS 2026 already announced, including which of this paper's mechanism-level guarantees still need to be confirmed under real field conditions rather than simulation.
